\documentclass{amsart}
% For dealing with references we use the comment environment
\usepackage{verbatim}
\newenvironment{reference}{\comment}{\endcomment}
%\newenvironment{reference}{}{}
\newenvironment{slogan}{\comment}{\endcomment}
\newenvironment{history}{\comment}{\endcomment}

% For commutative diagrams we use Xy-pic
\usepackage[all]{xy}

% We use 2cell for 2-commutative diagrams.
\xyoption{2cell}
\UseAllTwocells

% We use multicol for the list of chapters between chapters
\usepackage{multicol}

% This is generally recommended for better output
\usepackage{lmodern}
\usepackage[T1]{fontenc}

% For cross-file-references

% Package for hypertext links:
\usepackage{hyperref}

% For any local file, say "hello.tex" you want to link to please

% Theorem environments.
%
\theoremstyle{plain}
\newtheorem{theorem}[subsection]{Theorem}
\newtheorem{proposition}[subsection]{Proposition}
\newtheorem{lemma}[subsection]{Lemma}

\theoremstyle{definition}
\newtheorem{definition}[subsection]{Definition}
\newtheorem{example}[subsection]{Example}
\newtheorem{exercise}[subsection]{Exercise}
\newtheorem{situation}[subsection]{Situation}

\theoremstyle{remark}
\newtheorem{remark}[subsection]{Remark}
\newtheorem{remarks}[subsection]{Remarks}

\numberwithin{equation}{subsection}

% Macros
%
\def\lim{\mathop{\mathrm{lim}}\nolimits}
\def\colim{\mathop{\mathrm{colim}}\nolimits}
\def\Spec{\mathop{\mathrm{Spec}}}
\def\Hom{\mathop{\mathrm{Hom}}\nolimits}
\def\Ext{\mathop{\mathrm{Ext}}\nolimits}
\def\SheafHom{\mathop{\mathcal{H}\!\mathit{om}}\nolimits}
\def\SheafExt{\mathop{\mathcal{E}\!\mathit{xt}}\nolimits}
\def\Sch{\mathit{Sch}}
\def\Mor{\mathop{\mathrm{Mor}}\nolimits}
\def\Ob{\mathop{\mathrm{Ob}}\nolimits}
\def\Sh{\mathop{\mathit{Sh}}\nolimits}
\def\NL{\mathop{N\!L}\nolimits}
\def\CH{\mathop{\mathrm{CH}}\nolimits}
\def\proetale{{pro\text{-}\acute{e}tale}}
\def\etale{{\acute{e}tale}}
\def\QCoh{\mathit{QCoh}}
\def\Ker{\mathop{\mathrm{Ker}}}
\def\Im{\mathop{\mathrm{Im}}}
\def\Coker{\mathop{\mathrm{Coker}}}
\def\Coim{\mathop{\mathrm{Coim}}}

% Boxtimes
%
\DeclareMathSymbol{\boxtimes}{\mathbin}{AMSa}{"02}

%
% Macros for moduli stacks/spaces
%
\def\QCohstack{\mathcal{QC}\!\mathit{oh}}
\def\Cohstack{\mathcal{C}\!\mathit{oh}}
\def\Spacesstack{\mathcal{S}\!\mathit{paces}}
\def\Quotfunctor{\mathrm{Quot}}
\def\Hilbfunctor{\mathrm{Hilb}}
\def\Curvesstack{\mathcal{C}\!\mathit{urves}}
\def\Polarizedstack{\mathcal{P}\!\mathit{olarized}}
\def\Complexesstack{\mathcal{C}\!\mathit{omplexes}}
% \Pic is the operator that assigns to X its picard group, usage \Pic(X)
% \Picardstack_{X/B} denotes the Picard stack of X over B
% \Picardfunctor_{X/B} denotes the Picard functor of X over B
\def\Pic{\mathop{\mathrm{Pic}}\nolimits}
\def\Picardstack{\mathcal{P}\!\mathit{ic}}
\def\Picardfunctor{\mathrm{Pic}}
\def\Deformationcategory{\mathcal{D}\!\mathit{ef}}

\setlength{\emergencystretch}{3em}
\begin{document}
\title[Stacks Project, Tag 0C3V]{585 --- Integral closure of a reduced one-dimensional Nagata local ring commutes with henselization, strict henselization and completion}
\maketitle
\noindent Copyright (C) 2005--2025 Johan de Jong. Stacks Project authors. Source: \href{https://stacks.math.columbia.edu/tag/0C3V}{Tag 0C3V}.

\noindent Editorial difficulty note (not author text). Difficulty H1: The proof preserves the total fraction ring by a common denominator satisfying the normalization setup. It decomposes the finite base-changed normalization into local henselian or complete factors, proves them normal by identifying them with henselizations or completions of discrete valuation rings, and recovers the required integral closures..

\noindent Editorial provenance note (not author text). History: excerpt from revision \texttt{a0444\allowbreak 6e57e\allowbreak c1fbc\allowbreak 252a8\allowbreak 71afc\allowbreak ec775\allowbreak 2fb28\allowbreak 07b14}, varieties.tex, lines 8345--8423. Reference presentation was adapted; original-author context and core support blocks were retained or added. The mathematical wording is unchanged. Licensed under GNU FDL 1.2 or later; no invariant sections or cover texts. See ../licenses/COPYING. Context blocks reproduce author definitions and supporting results; they are not additional counted problems.

\begin{lemma}
\label{lemma-pre-delta-invariant}
Let $(A, \mathfrak m)$ be a reduced Nagata $1$-dimensional local ring.
Let $A'$ be the integral closure of $A$ in the total ring of fractions
of $A$. Then $A'$ is a normal Nagata ring, $A \to A'$ is finite, and
$A'/A$ has finite length as an $A$-module.
\end{lemma}

\begin{lemma}
\label{lemma-normalization-same-after-completion}
Let $A$ be a reduced Nagata local ring of dimension $1$.
Let $A \to A'$ be as in Lemma \ref{lemma-pre-delta-invariant}.
Let $A^h$, $A^{sh}$, resp.\ $A^\wedge$
be the henselization, strict henselization, resp.\ completion of $A$.
Then $A^h$, $A^{sh}$, resp. $A^\wedge$ is a reduced Nagata local
ring of dimension $1$ and
$A' \otimes_A A^h$, $A' \otimes_A A^{sh}$, resp. $A' \otimes_A A^\wedge$
is the integral closure of $A^h$, $A^{sh}$, resp.\ $A^\wedge$
in its total ring of fractions.
\end{lemma}

\begin{proof}
Observe that $A^\wedge$ is reduced, see
More on Algebra, Lemma \href{https://stacks.math.columbia.edu/tag/07NZ}{lemma completion reduced (Tag 07NZ)}.
The rings $A^h$ and $A^{sh}$ are reduced by
More on Algebra, Lemma \href{https://stacks.math.columbia.edu/tag/06DH}{lemma henselization reduced (Tag 06DH)}.
The dimensions of $A$, $A^h$, $A^{sh}$, and $A^\wedge$ are the same
by More on Algebra, Lemmas
\href{https://stacks.math.columbia.edu/tag/07NV}{lemma completion dimension (Tag 07NV)} and
\href{https://stacks.math.columbia.edu/tag/06LK}{lemma henselization dimension (Tag 06LK)}.

\medskip\noindent
Recall that a Noetherian local ring is Nagata if and only if the
formal fibres of $A$ are geometrically reduced, see
More on Algebra, Lemma \href{https://stacks.math.columbia.edu/tag/0BJ0}{lemma Nagata local ring (Tag 0BJ0)}.
This property is inherited by $A^h$ and $A^{sh}$, see
the material in More on Algebra, Section
\href{https://stacks.math.columbia.edu/tag/0BIR}{section properties formal fibres (Tag 0BIR)}
and especially Lemma \href{https://stacks.math.columbia.edu/tag/0C36}{lemma henselization P ring (Tag 0C36)}.
The completion is Nagata by
Algebra, Lemma \href{https://stacks.math.columbia.edu/tag/032W}{lemma Noetherian complete local Nagata (Tag 032W)}.

\medskip\noindent
Now we come to the statement on integral closures. Before continuing
let us pick $f \in \mathfrak m$ as in
Lemma \href{https://stacks.math.columbia.edu/tag/0C3R}{lemma pre pre delta invariant (Tag 0C3R)}.
Then the image of $f$ in $A^h$, $A^{sh}$, and $A^\wedge$
clearly is an element satisfying properties (1) -- (6)
in that ring.

\medskip\noindent
Since $A \to A'$ is finite we see that
$A' \otimes_A A^h$ and $A' \otimes_A A^{sh}$
is the product of henselian local rings finite over $A^h$ and
$A^{sh}$, see Algebra, Lemma \href{https://stacks.math.columbia.edu/tag/04GH}{lemma finite over henselian (Tag 04GH)}.
Each of these local rings is the henselization of $A'$ at a
maximal ideal $\mathfrak m' \subset A'$ lying over $\mathfrak m$, see
Algebra, Lemma \href{https://stacks.math.columbia.edu/tag/05WP}{lemma quasi finite henselization (Tag 05WP)} or
\href{https://stacks.math.columbia.edu/tag/05WR}{lemma quasi finite strict henselization (Tag 05WR)}.
Hence these local rings are normal domains by
More on Algebra, Lemma \href{https://stacks.math.columbia.edu/tag/06DI}{lemma henselization normal (Tag 06DI)}.
It follows that $A' \otimes_A A^h$ and $A' \otimes_A A^{sh}$
are normal rings. Since $A^h \to A' \otimes_A A^h$ and
$A^{sh} \to A' \otimes_A A^{sh}$ are finite (hence integral) and since
$A' \otimes_A A^h \subset (A^h)_f = Q(A^h)$ and
$A' \otimes_A A^{sh} \subset (A^{sh})_f = Q(A^{sh})$
we conclude that $A' \otimes_A A^h$ and $A' \otimes_A A^{sh}$ are
the desired integral closures.

\medskip\noindent
For the completion we argue in entirely the same manner. First,
by Algebra, Lemma \href{https://stacks.math.columbia.edu/tag/07N9}{lemma completion finite extension (Tag 07N9)} we have
$$
A' \otimes_A A^\wedge = (A')^\wedge =
\prod\nolimits (A'_{\mathfrak m'})^\wedge
$$
The local rings $A'_{\mathfrak m'}$ are normal and have dimension $1$
(by Algebra, Lemma \href{https://stacks.math.columbia.edu/tag/02MA}{lemma finite in codim 1 (Tag 02MA)} for example or
the discussion in
Algebra, Section \href{https://stacks.math.columbia.edu/tag/00OG}{section homomorphism dimension (Tag 00OG)}).
Thus $A'_{\mathfrak m'}$ is a discrete valuation ring, see
Algebra, Lemma \href{https://stacks.math.columbia.edu/tag/00PD}{lemma characterize dvr (Tag 00PD)}.
Hence $(A'_{\mathfrak m'})^\wedge$ is a discrete valuation ring
by More on Algebra, Lemma \href{https://stacks.math.columbia.edu/tag/0AP1}{lemma completion dvr (Tag 0AP1)}.
It follows that $A' \otimes_A A^\wedge$ is a normal ring and
we can conclude in exactly the same manner as before.
\end{proof}
\end{document}
